\pagestyle{empty}
\tight
\begin{tblr}{width=\textwidth,colspec={|Q[c,m,1.9cm]|X[l,m]|},hlines,vlines,hspan=minimal,row{1}={bg=headblue},rowsep=1pt,colsep=3pt}
\SetCell{c,m}\hfil\logo\hfil & \textbf{POLITEKNIK NEGERI MALANG}\par\textbf{JURUSAN TEKNOLOGI INFORMASI}\par\textbf{PROGRAM STUDI: D4 TEKNIK INFORMATIKA} \\
\SetCell[c=2]{c} \textbf{RENCANA PEMBELAJARAN SEMESTER (RPS)} & \\
\end{tblr}
\begin{longtblr}{width=\textwidth,colspec={|Q[l,m,1.9cm]|X[0.85,l,m]|X[1.45,l,m]|X[0.75,l,m]|X[1.15,l,m]|},hlines,vlines,hspan=minimal,rowsep=1pt,colsep=2pt,row{1,3}={bg=headgray,font=\bfseries}}
MATA KULIAH & KODE & BOBOT (SKS)/jam & SEMESTER & TGL. PENYUSUNAN \\
Pancasila & RTI251001 & 2 SKS/2 jam & 1 & 11 Agustus 2025 \\
OTORISASI & Dosen Pengembang RPS & Koordinator RMK & \SetCell[c=2]{l} Ka PRODI &  \\
 & Dr. Widaningsih Condrowardani, S.Psi., S.H., M.H. & Dr. Widaningsih Condrowardani, S.Psi., S.H., M.H. & \SetCell[c=2]{l}{\vspace{8mm}\par Dr. Ely Setyo Astuti, S.T., M.T.} &  \\
\SetCell[r=9]{l} Capaian Pembelajaran (CP) & \SetCell[c=4]{l,bg=headgray,font=\bfseries} Capaian Pembelajaran Lulusan Program Studi (CPL-Prodi) &  &  &  \\
 & CPL01 & \SetCell[c=3]{l} Mampu menginternalisasi nilai ketakwaan, ketaatan hukum, kedisiplinan, profesionalisme, etika profesi, semangat belajar sepanjang hayat, serta kepedulian terhadap isu sosial dan dampak teknologi. &  &  \\
 & \SetCell[c=4]{l,bg=headgray,font=\bfseries} Capaian Pembelajaran mata kuliah (CPL-MK) &  &  &  \\
 & CPMK0102 & \SetCell[c=3]{l} Mahasiswa mampu menginternalisasi nilai ketakwaan, Pancasila, dan kewarganegaraan serta menunjukkan semangat belajar sepanjang hayat, tanggung jawab sosial, dan kemampuan adaptasi karir. &  &  \\
 & \SetCell[c=4]{l,bg=headgray,font=\bfseries} Kemampuan akhir tiap tahapan belajar (Sub-CPMK) &  &  &  \\
 & SCPMK0102-00101 & \SetCell[c=3]{l} Mahasiswa mampu menjelaskan Pendidikan Pancasila dalam tinjauan historis, kultural, yuridis, dan filosofis, konteks sejarah perjuangan bangsa Indonesia, serta Pancasila sebagai sistem filsafat [C2, A3]. &  &  \\
 & SCPMK0102-00102 & \SetCell[c=3]{l} Mahasiswa mampu menjelaskan UUD RI 1945, amandemen UUD RI 1945, Trias Politika, dan kelembagaan negara menurut UUD RI 1945 [C2, A3]. &  &  \\
 & SCPMK0102-00103 & \SetCell[c=3]{l} Mahasiswa mampu menjelaskan Pancasila sebagai ideologi nasional dan membandingkannya dengan ideologi lain yang berkembang di dunia [C2, A3]. &  &  \\
 & SCPMK0102-00104 & \SetCell[c=3]{l} Mahasiswa mampu menginternalisasi nilai Pancasila dalam isu HAM, pencegahan tindak pidana korupsi, dan Pancasila sebagai paradigma pembangunan [C3, A4]. &  &  \\
\SetCell[r=2]{l} Penilaian & \SetCell[c=4]{l} *Nilai CPMK didapat dari agregasi nilai SUB-CPMK yang dibebankan &  &  &  \\
 & \SetCell[c=4]{l}{\tiny\begin{tblr}{width=\linewidth,colspec={|Q[c,m,0.1\linewidth]|Q[c,m,0.16\linewidth]|X[l,m]|X[l,m]|X[l,m]|X[l,m]|X[l,m]|Q[c,m,0.08\linewidth]|},hlines,vlines,hspan=minimal,rowsep=1pt,colsep=0.7pt,row{1,2}={bg=headgray,font=\bfseries}}
Kode CPMK & Kode SCPMK & \SetCell[c=5]{c} Bobot per Bentuk Penilaian & & & & & Total Bobot \\
 & & Tugas & Kuis & Case Method & UTS & UAS & \\
\SetCell[r=1]{c} \SetCell[r=4]{c} CPMK0102 & SCPMK0102-00101 & 5\%: kajian historis-kultural-yuridis-filosofis & 5\%: konsep dasar dan filsafat Pancasila & & 10\%: materi minggu 1--7 & & 20\% \\
\SetCell[r=3]{c}  & SCPMK0102-00102 & 10\%: kajian UUD 1945 dan kelembagaan negara & & & 10\%: materi minggu 1--7 & & 20\% \\
 & SCPMK0102-00103 & 2.5\%: kajian ideologi nasional & 5\%: ideologi nasional dan ideologi dunia & & & 15\%: materi minggu 9--16 & 22.5\% \\
 & SCPMK0102-00104 & 2.5\%: kajian Pancasila dan HAM & & 15\%: kasus HAM, korupsi, dan pembangunan & & 20\%: materi minggu 9--16 & 37.5\% \\
\SetCell[c=7]{r}\textbf{Total Per Penilaian} & & & & & & & \textbf{100\%} \\
\end{tblr}} &  &  &  \\
Diskripsi Singkat Mata Kuliah & \SetCell[c=4]{l} Mata kuliah Pancasila membekali mahasiswa untuk memahami sejarah, landasan, filosofi, dan kedudukan Pancasila sebagai dasar Negara Republik Indonesia serta menginternalisasi nilai-nilainya dalam isu HAM, antikorupsi, dan paradigma pembangunan. &  &  &  \\
Bahan Kajian / Materi Pembelajaran & \SetCell[c=4]{l}\begin{enumerate}\item Pendidikan Pancasila dalam tinjauan historis, kultural, yuridis, dan filosofis\item Pancasila dalam konteks sejarah perjuangan bangsa Indonesia\item Pancasila sebagai sistem filsafat\item UUD RI 1945 dan amandemen UUD RI 1945\item Trias Politika dan kelembagaan negara menurut UUD RI 1945\item Pancasila sebagai ideologi nasional dan ideologi lain yang berkembang di dunia\item Pancasila dan HAM serta pelaksanaan HAM dalam UUD RI 1945\item Tindak pidana korupsi\item Pancasila sebagai paradigma pembangunan\end{enumerate} &  &  &  \\
\SetCell[r=2]{l} Pustaka & \SetCell[c=4]{l} \textbf{Utama:}\begin{enumerate}\item Modul Ajar Pendidikan Pancasila.\end{enumerate} &  &  &  \\
 & \SetCell[c=4]{l} \textbf{Pendukung:}\begin{enumerate}\item Kementerian Riset, Teknologi, dan Pendidikan Tinggi Republik Indonesia. 2016. Buku Ajar Pendidikan Pancasila untuk Perguruan Tinggi. Jakarta: Direktorat Pembelajaran dan Kemahasiswaan Kemenristekdikti.\end{enumerate} &  &  &  \\
Media Pembelajaran & \SetCell[c=2]{l} \textbf{Software:} Microsoft Word, Microsoft PowerPoint, LMS. &  & \SetCell[c=2]{l} \textbf{Hardware:} Komputer/laptop, proyektor. &  \\
Nama Dosen Pengampu & \SetCell[c=4]{l}\begin{enumerate}\item Dr. Widaningsih Condrowardani, S.Psi., S.H., M.H.\end{enumerate} &  &  &  \\
Matakuliah Syarat & \SetCell[c=4]{l} - &  &  &  \\
\end{longtblr}
\begin{longtblr}{width=\textwidth,colspec={|Q[c,m,0.06\textwidth]|X[1.35,l,m]|X[1.08,l,m]|X[1.16,l,m]|X[0.7,l,m]|X[1.24,l,m]|X[1.06,l,m]|X[0.98,l,m]|Q[c,m,0.05\textwidth]|},hlines,vlines,hspan=minimal,rowhead=2,row{1,2}={bg=headgreen,font=\bfseries},rowsep=1pt,colsep=1.5pt}
Mgg.\par Ke & Kemampuan Akhir Yang Direncanakan (Sub-CP-MK) & Bahan kajian (Materi Pembelajaran) & Bentuk dan Metode Pembelajaran & Estimasi Waktu & Pengalaman Belajar Mahasiswa & Kriteria \& Bentuk Penilaian & Indikator Penilaian & Bobot Penilaian (\%) \\
(1) & (2) & (3) & (4) & (5) & (6) & (7) & (8) & (9) \\
1 & SCPMK0102-00101 & Tujuan akhir perkuliahan; definisi Pancasila secara historis, kultural, yuridis, dan filosofis. & Bentuk: Kuliah/Luring. Metode: Ceramah, presentasi, diskusi, dan tanya jawab. & TM: 1 x (2 x 50'); BT+BM: 1 x (2 x 60'). & Mahasiswa mengenal tujuan akhir perkuliahan serta menjelaskan definisi Pancasila secara historis, kultural, yuridis, dan filosofis. & Bentuk: tes lisan dan tugas. Kriteria: ketepatan dan kejelasan penjelasan. Mengacu rubrik RTM. & Pemahaman definisi Pancasila secara historis, kultural, yuridis, dan filosofis; kejelasan menjawab pertanyaan. & 2.5 \\
2 & SCPMK0102-00101 & Pancasila dalam konteks sejarah perjuangan bangsa Indonesia; prinsip-prinsip Pancasila. & Bentuk: Kuliah/Luring. Metode: Ceramah, presentasi, diskusi, dan tanya jawab. & TM: 1 x (2 x 50'); BT+BM: 1 x (2 x 60'). & Mahasiswa memahami Pancasila dalam konteks sejarah perjuangan bangsa Indonesia serta prinsip-prinsip Pancasila. & Bentuk: tes lisan dan tugas. Kriteria: ketepatan penjelasan, kreativitas, dan daya tarik penyampaian. Mengacu rubrik RTM. & Kemampuan memahami Pancasila dalam konteks sejarah perjuangan bangsa dan prinsip-prinsip Pancasila. & 2.5 \\
3 & SCPMK0102-00101 & Pancasila sebagai sistem filsafat; unsur filsafat, hierarki nilai, dan penerapannya sehari-hari. Kuis 1. & Bentuk: Kuliah/Luring. Metode: Ceramah, presentasi, diskusi, tanya jawab, dan kuis. & TM: 1 x (2 x 50'); BT+BM: 1 x (2 x 60'). & Mahasiswa memahami Pancasila sebagai sistem filsafat, hierarki nilai, dan penerapannya dalam kehidupan sehari-hari, lalu mengerjakan kuis. & Bentuk: kuis tertulis. Kriteria: ketepatan jawaban dan pemecahan masalah. Mengacu rubrik RTM. & Kemampuan memahami unsur filsafat dalam Pancasila, hierarki nilai, dan penerapannya. & 5 \\
4 & SCPMK0102-00102 & UUD RI 1945: sejarah perumusan dan pengesahan, pokok pikiran pembukaan, nilai, dan kedudukan. & Bentuk: Kuliah/Luring. Metode: Ceramah, presentasi, diskusi, dan tanya jawab. & TM: 1 x (2 x 50'); BT+BM: 1 x (2 x 60'). & Mahasiswa memahami definisi, sejarah perumusan, pokok-pokok pikiran pembukaan, nilai, dan kedudukan UUD 1945. & Bentuk: tes tulis/lisan dan tugas. Kriteria: kejelasan jawaban dan kesesuaian isi tulisan dengan materi. Mengacu rubrik RTM. & Kemampuan memahami definisi, sejarah, pokok pikiran pembukaan, nilai, dan kedudukan UUD 1945. & 2.5 \\
5 & SCPMK0102-00102 & Amandemen UUD RI 1945: definisi, sejarah, tujuan, dan proses amandemen. & Bentuk: Kuliah/Luring. Metode: Ceramah, presentasi, diskusi, dan tanya jawab. & TM: 1 x (2 x 50'); BT+BM: 1 x (2 x 60'). & Mahasiswa mengetahui definisi, sejarah, tujuan, dan proses amandemen UUD 1945. & Bentuk: tes lisan dan tugas. Kriteria: kemampuan melakukan presentasi. Mengacu rubrik RTM. & Kemampuan memahami definisi, sejarah, tujuan, dan proses amandemen UUD 1945. & 2.5 \\
6 & SCPMK0102-00102 & Trias Politika dalam Negara RI: definisi, penerapan di Indonesia, dan urgensinya. & Bentuk: Kuliah/Luring. Metode: Ceramah, presentasi, diskusi, dan tanya jawab. & TM: 1 x (2 x 50'); BT+BM: 1 x (2 x 60'). & Mahasiswa mengetahui definisi Trias Politika, penerapannya di Indonesia, dan pentingnya memahami Trias Politika. & Bentuk: tes lisan dan tugas. Kriteria: ketepatan dan kejelasan menjawab pertanyaan. Mengacu rubrik RTM. & Kemampuan memahami definisi, penerapan, dan urgensi Trias Politika dalam Negara RI. & 2.5 \\
7 & SCPMK0102-00102 & Kelembagaan negara menurut UUD RI 1945 dan lembaga-lembaga negara utama. & Bentuk: Kuliah/Luring. Metode: Ceramah, presentasi, diskusi, dan tanya jawab. & TM: 1 x (2 x 50'); BT+BM: 1 x (2 x 60'). & Mahasiswa mengetahui definisi kelembagaan negara menurut UUD RI 1945 dan lembaga-lembaga negara utama. & Bentuk: tes lisan dan tugas. Kriteria: kelancaran dan kesesuaian dalam kerja tim. Mengacu rubrik RTM. & Kemampuan memahami kelembagaan negara menurut UUD RI 1945 dan lembaga negara utama. & 2.5 \\
8 & SCPMK0102-00101, SCPMK0102-00102 & UTS: Pendidikan Pancasila tinjauan historis-kultural-yuridis-filosofis, sejarah perjuangan bangsa, sistem filsafat, UUD RI 1945, amandemen, Trias Politika, dan kelembagaan negara. & Bentuk: Ujian tertulis (esai dan/atau pilihan ganda). & TM: 1 x (2 x 50'). & Mahasiswa menjelaskan dan menjawab soal ujian tertulis esai/pilihan ganda terkait materi minggu 1--7. & Bentuk: tes tulis pilihan ganda/esai. Kriteria: ketepatan jawaban. Mengacu rubrik RTM. & Kemampuan menjawab soal ujian tertulis esai/pilihan ganda terkait materi minggu 1--7 dengan benar. & 20 \\
9 & SCPMK0102-00103 & Pancasila sebagai ideologi nasional: definisi, fungsi, dan proses terbentuknya. & Bentuk: Kuliah/Luring. Metode: Ceramah, presentasi, diskusi, dan tanya jawab. & TM: 1 x (2 x 50'); BT+BM: 1 x (2 x 60'). & Mahasiswa mengetahui definisi, fungsi, dan proses terbentuknya Pancasila sebagai ideologi nasional. & Bentuk: tes lisan dan tugas. Kriteria: kelancaran dan kesesuaian dalam kerja tim. Mengacu rubrik RTM. & Kemampuan memahami definisi, fungsi, dan proses terbentuknya Pancasila sebagai ideologi nasional. & 2.5 \\
10 & SCPMK0102-00103 & Ideologi lain yang berkembang di dunia: definisi, urgensi mempelajari ideologi, dan macam-macam ideologi. Kuis 2. & Bentuk: Kuliah/Luring. Metode: Ceramah, presentasi, diskusi, tanya jawab, dan kuis. & TM: 1 x (2 x 50'); BT+BM: 1 x (2 x 60'). & Mahasiswa mengetahui macam-macam ideologi dunia, membandingkannya dengan Pancasila, lalu mengerjakan kuis. & Bentuk: kuis tertulis. Kriteria: ketepatan jawaban. Mengacu rubrik RTM. & Kemampuan memahami macam-macam ideologi dunia dan membandingkannya dengan ideologi Pancasila. & 5 \\
11 & SCPMK0102-00104 & Pancasila dan HAM: definisi, hubungan, implementasi, bentuk-bentuk HAM, dan contoh penerapannya. & Bentuk: Kuliah/Luring. Metode: Ceramah, presentasi, diskusi, dan tanya jawab. & TM: 1 x (2 x 50'); BT+BM: 1 x (2 x 60'). & Mahasiswa mengetahui hubungan Pancasila dan HAM, bentuk-bentuk HAM, serta contoh penerapan HAM. & Bentuk: tes lisan dan tugas. Kriteria: kelancaran dan kesesuaian dalam kerja tim. Mengacu rubrik RTM. & Kemampuan memahami hubungan Pancasila dan HAM, bentuk HAM, dan contoh penerapannya. & 2.5 \\
12 & SCPMK0102-00104 & Pelaksanaan HAM dalam UUD RI 1945: jaminan HAM dalam pasal UUD 1945 dan pengawasan pelaksanaannya. & Bentuk: Kuliah/Luring. Metode: Case Method, FGD, presentasi, dan diskusi kelompok. & TM: 1 x (2 x 50'); BT+BM: 1 x (2 x 60'). & Mahasiswa menganalisis kasus pelaksanaan HAM berdasarkan jaminan HAM dalam pasal-pasal UUD 1945. & Bentuk: penilaian case method. Kriteria: ketajaman analisis kasus dan argumentasi. Mengacu rubrik RTM. & Kemampuan menganalisis jaminan HAM dalam pasal UUD 1945 dan pengawasan pelaksanaan HAM di Indonesia. & 3 \\
13 & SCPMK0102-00104 & Tindak pidana korupsi: pengertian, jenis-jenis, dan dasar hukum. & Bentuk: Kuliah/Luring. Metode: Case Method, FGD, presentasi, dan diskusi kelompok. & TM: 1 x (2 x 50'); BT+BM: 1 x (2 x 60'). & Mahasiswa menganalisis kasus tindak pidana korupsi berdasarkan pengertian, jenis, dan dasar hukumnya. & Bentuk: penilaian case method. Kriteria: kelancaran dan kesesuaian dalam kerja tim. Mengacu rubrik RTM. & Kemampuan memahami pengertian, jenis-jenis, dan dasar hukum tindak pidana korupsi. & 3 \\
14 & SCPMK0102-00104 & Upaya pencegahan serta penanggulangan tindak pidana korupsi. & Bentuk: Kuliah/Luring. Metode: Case Method, FGD, presentasi, dan diskusi kelompok. & TM: 1 x (2 x 50'); BT+BM: 1 x (2 x 60'). & Mahasiswa merumuskan upaya pencegahan dan penanggulangan korupsi dari kasus aktual berbasis nilai Pancasila. & Bentuk: penilaian case method. Kriteria: kelancaran dan kesesuaian dalam kerja tim. Mengacu rubrik RTM. & Kemampuan merumuskan upaya pencegahan serta penanggulangan tindak pidana korupsi. & 3 \\
15 & SCPMK0102-00104 & Pancasila sebagai paradigma pembangunan: definisi dan peranan Pancasila dalam pembangunan di Indonesia. & Bentuk: Kuliah/Luring. Metode: Case Method, FGD, presentasi, dan diskusi kelompok. & TM: 1 x (2 x 50'); BT+BM: 1 x (2 x 60'). & Mahasiswa memahami definisi dan peranan Pancasila sebagai paradigma pembangunan di Indonesia. & Bentuk: penilaian case method. Kriteria: kelancaran dan kesesuaian dalam kerja tim. Mengacu rubrik RTM. & Kemampuan memahami definisi dan peranan Pancasila sebagai paradigma pembangunan. & 3 \\
16 & SCPMK0102-00104 & Contoh paradigma pembangunan dan presentasi hasil analisis kasus. & Bentuk: Kuliah/Luring. Metode: Case Method, FGD, presentasi, dan diskusi kelompok. & TM: 1 x (2 x 50'); BT+BM: 1 x (2 x 60'). & Mahasiswa mempresentasikan contoh paradigma pembangunan dan hasil analisis kasus berbasis nilai Pancasila. & Bentuk: penilaian case method. Kriteria: daya tarik presentasi dan kesesuaian materi. Mengacu rubrik RTM. & Kemampuan menjelaskan contoh paradigma pembangunan dan mempresentasikan hasil analisis kasus. & 3 \\
17 & SCPMK0102-00103, SCPMK0102-00104 & UAS: ideologi nasional, ideologi lain di dunia, Pancasila dan HAM, pelaksanaan HAM dalam UUD RI 1945, tindak pidana korupsi, dan Pancasila sebagai paradigma pembangunan. & Bentuk: Ujian tertulis (esai dan/atau pilihan ganda). & TM: 1 x (2 x 50'). & Mahasiswa menjelaskan dan menjawab soal ujian tertulis esai/pilihan ganda terkait materi minggu 9--16. & Bentuk: tes tulis pilihan ganda/esai. Kriteria: jawaban benar dan logis. Mengacu rubrik RTM. & Kemampuan menjawab pertanyaan ujian dengan benar dan logis terkait materi minggu 9--16. & 35 \\
\end{longtblr}
